\documentclass[../../main.tex]{subfiles}% 注意这里的文件路径不能用 ./main.tex ,否则用latexmk编译子文件会报错
\graphicspath{{\subfix{./image/}}{\subfix{../../image/}}} % 指定图片引用路径,后续可以直接使用图片文件名
% 如果图片存放在上级目录,则使用 ../../image/ 作为备用路径

% 例如:
% \begin{figure}[H]
% \centering
% \includegraphics[width=0.3\textwidth]{图.png}
% \caption{}
% \label{figure:图}
% \end{figure}
% 注意:上述\label{}一定要放在\caption{}之后,否则引用图片序号会只会显示??.

\begin{document}

\section{度量外测度与容量}\label{section:ch01-度量外测度容量}

\subsection{容量}\label{subsection:ch01-容量}

\begin{definition}\label{definition:ch01-容量}
设 $(X,d)$ 为一度量空间，$\mathcal{F}\subseteq\mathcal{P}(X)$ 满足以下性质：(a) 紧子集族 $\mathcal{K}$ 包含于 $\mathcal{F}$；(b) 对每个序列 $\{E_n\}\subseteq\mathcal{F}$，有 $\bigcup_{n=1}^\infty E_n\in\mathcal{F}$。映射 $C:\mathcal{F}\to[0,\infty]$ 称为 $\mathcal{F}$ 上的一个容量，若：
\begin{enumerate}[(i)]
\item $C(\varnothing )=0$；
\item 对任意 $E_1,E_2\in\mathcal{F}$、$E_1\subseteq E_2$，有 $C(E_1)\leqslant C(E_2)$；
\item 对任意序列 $\{E_n\}\subseteq\mathcal{F}$，有
\[
C\left(\bigcup_{n=1}^\infty E_n\right)\leqslant\sum_{n=1}^\infty C(E_n).
\]
\end{enumerate}
$\mathcal{P}(X)$ 上的容量称为 $X$ 上的容量。
\end{definition}

\begin{remark}\label{remark:ch01-容量与外测度}
若 $\mathcal{F}$ 满足上述性质 (a)--(b)，则由附录 A 知开子集族 $\mathcal{T}$ 包含于 $\mathcal{F}$。此外，由\cref{definition:ch01-外测度} 和\cref{definition:ch01-容量}，$X$ 上的容量是一个外测度，而度量空间 $X$ 上的外测度是一个容量。
\end{remark}

\begin{definition}\label{definition:ch01-内外容量}
一个容量称为外容量，若对任意 $E\in\mathcal{F}$，
\[
C(E)=\inf\{C(A)\mid A\supseteq E,\ A\in\mathcal{T}\};
\]
称为内容量，若对任意 $E\in\mathcal{F}$，
\[
C(E)=\sup\{C(K)\mid K\subseteq E,\ K\in\mathcal{K}\}.
\]
集合 $E\in\mathcal{F}$ 称为 $C$--可容的，若
\begin{align}
C(E)=\inf\{C(A)\mid A\supseteq E,\ A\in\mathcal{T}\}=\sup\{C(K)\mid K\subseteq E,\ K\in\mathcal{K}\}. \label{align:ch01-可容}
\end{align}
\end{definition}

\begin{remark}\label{remark:ch01-容量生成}
设 $X$ 和 $\mathcal{F}\subseteq\mathcal{P}(X)$ 如\cref{definition:ch01-容量} 所述，$C$ 为 $\mathcal{F}$ 上的一个容量。对任意族 $\mathcal{F}_1\subseteq\mathcal{F}$，定义 $C_{\mathcal{F}_1}:\mathcal{P}(X)\to[0,\infty]$ 为
\[
C_{\mathcal{F}_1}(E)\triangleq \inf\{C(F)\mid F\supseteq E,\ F\in\mathcal{F}_1\}\qquad(E\subseteq X)
\]
若存在 $F\in\mathcal{F}_1$ 使得 $F\supseteq E$，否则令 $C_{\mathcal{F}_1}(E)\triangleq \infty$。由\cref{theorem:ch01-外测度生成}，$C_{\mathcal{F}_1}$ 是一个外测度，因而是 $X$ 上的一个容量（见\cref{remark:ch01-容量与外测度}）。特别地，若 $\mathcal{F}_1=\mathcal{T}$ 为 $X$ 上的拓扑，则 $C_{\mathcal{T}}$ 由定义是一个外容量，且 $C$ 是外容量当且仅当 $C=C_{\mathcal{T}}$。

\cref{definition:ch01-内外容量} 中的概念与\cref{definition:ch01-Borel测度}(b)--(d) 中的概念类似。
\end{remark}

在第 3 章中，我们将使用以下结果（见文献~[2, Theorem 2.3.11]、[32]、[109] 及其中相关引文）。

\begin{theorem}[Choquet]\label{theorem:ch01-Choquet}
设 $C$ 为 $\mathbb{R}^N$ 上的一个容量，满足：
\begin{enumerate}[(i)]
\item 对任意非增紧子集序列 $\{K_k\}\subseteq\mathbb{R}^N$，
\begin{align}
C\left(\bigcap_{k=1}^\infty K_k\right)=\lim_{k\to\infty}C(K_k); \label{align:ch01-Choquet上连续}
\end{align}
\item 对任意非减序列 $\{E_k\}\subseteq\mathcal{P}(\mathbb{R}^N)$，
\begin{align}
C\left(\bigcup_{k=1}^\infty E_k\right)=\lim_{k\to\infty}C(E_k). \label{align:ch01-Choquet下连续}
\end{align}
\end{enumerate}
则每个 $E\in\mathcal{B}(\mathbb{R}^N)$ 都是 $C$--可容的。
\end{theorem}

\begin{remark}\label{remark:ch01-Suslin集}
\cref{theorem:ch01-Choquet} 的结论实际上对更大的 Suslin 集类成立。
\end{remark}

关于\cref{theorem:ch01-Choquet}，我们注意以下结果。

\begin{proposition}\label{proposition:ch01-外容量上连续}\label{proposition:ch01-测度上连续}
设 $X$ 和 $\mathcal{F}\subseteq\mathcal{P}(X)$ 如\cref{definition:ch01-容量} 所述，$C:\mathcal{F}\to[0,\infty]$ 为一外容量。则对任意非增紧集序列 $\{K_n\}\subseteq\mathcal{K}$，等式\eqref{align:ch01-Choquet上连续} 成立。
\end{proposition}

\begin{proof}
令 $K\triangleq \bigcap_{n=1}^\infty K_n$，考虑任意 $A\supseteq K$、$A\in\mathcal{T}$。由于 $\{K_n\}$ 非增，存在 $\bar{n}\in\mathbb{N}$ 使得对所有 $n>\bar{n}$ 有 $K_n\subseteq A$，故由 $C$ 的单调性显然有
\[
C(K)\leqslant\lim_{n\to\infty}C(K_n)\leqslant\inf_{A\supseteq K,\ A\in\mathcal{T}}C(A).
\]
另一方面，由假设有 $\inf_{A\supseteq K,\ A\in\mathcal{T}}C(A)=C(K)$。于是结论成立。
\end{proof}

设 $X$ 和 $\mathcal{F}\subseteq\mathcal{P}(X)$ 如\cref{definition:ch01-容量} 所述，$C:\mathcal{F}\to[0,\infty]$ 为一容量。以下概念与\cref{subsection:ch01-零测集} 中的概念类似：
\begin{enumerate}[(a)]
\item 满足 $C(E)=0$ 的集合 $E\in\mathcal{F}$ 称为 $C$--零测的；
\item 性质 $P$ 称为 $C$--拟处处成立（记为 $C$--q.e. 或 q.e.），若集合 $\{x\in U\mid P(x)\text{ 不成立}\}$ 是 $C$--零测的。
\end{enumerate}

由\cref{definition:ch01-容量}，$C$--零测集族 $\mathcal{N}_C$ 关于可数并和遗传性是稳定的，即 $F\subseteq E$、$E\in\mathcal{N}_C\Rightarrow F\in\mathcal{N}_C$。(b) 的特殊情形如下：
\begin{itemize}
\item 设 $f,f_n:X\to\mathbb{R}$（$n\in\mathbb{N}$）。我们说 $f_n$ 在 $X$ 中 $C$--拟处处收敛于 $f$，若存在 $N\in\mathcal{N}_C$ 使得对任意 $x\in N^c$ 有 $f_n(x)$ 收敛于 $f(x)$；
\item 两个函数 $f,g:X\to\mathbb{R}$ 在 $X$ 中 $C$--q.e. 相等，若 $\{x\in X\mid f(x)\neq g(x)\}\in\mathcal{N}_C$。
\end{itemize}

显然，$C$--q.e. 相等是一个等价关系：
\begin{align}
g\sim f\overset{\mathrm{def}}{\Longleftrightarrow}g=f\ C\text{--q.e. 在 }X\text{ 中}, \label{align:ch01-q.e.等价}
\end{align}
关于\eqref{align:ch01-q.e.等价} 的每个等价类由其任意一个元素唯一确定，即该类的一个代表元。与 $\mu$--a.e. 相等的函数一样，我们将 $C$--q.e. 相等的函数等同起来，从而将整个类视为在 $X$ 中 $C$--q.e. 定义的唯一映射 $f$。

\subsection{度量外测度}\label{subsection:ch01-度量外测度}

\begin{definition}\label{definition:ch01-度量外测度}
设 $(X,d)$ 为一度量空间。外测度 $\mu^*$ 称为度量外测度，若对任意非空集合 $E,F\subseteq X$ 且 $d(E,F)>0$，
\begin{align}
\mu^*(E\cup F)=\mu^*(E)+\mu^*(F). \label{align:ch01-度量外测度条件}
\end{align}
\end{definition}

度量外测度的存在性可如\cref{subsection:ch01-Caratheodory构造} 中那样，通过对定义\eqref{align:ch01-外测度生成a} 作微小修改来证明。设 $(X,d)$ 为一度量空间，$\mathcal{C}\subseteq\mathcal{P}(X)$ 且 $\varnothing \in\mathcal{C}$，$\eta:\mathcal{C}\to[0,\infty]$ 且 $\eta(\varnothing )=0$。对任意 $\delta\in(0,\infty]$，令
\begin{align}
\eta_\delta^*(E)\triangleq \inf\left\{\sum_{n=1}^\infty\eta(E_n)\ \middle|\ E\subseteq\bigcup_{n=1}^\infty E_n,\ \{E_n\}\subseteq\mathcal{C},\ \operatorname{diam}(E_n)\leqslant\delta\ \forall n\in\mathbb{N}\right\} \label{align:ch01-eta-delta}
\end{align}
若存在如\eqref{align:ch01-eta-delta} 中 $E$ 的可数覆盖 $\{E_n\}$，否则令 $\eta_\delta^*(E)\triangleq \infty$。如\cref{theorem:ch01-外测度生成} 的证明那样论证，可知映射 $\eta_\delta^*:\mathcal{P}(X)\to[0,\infty]$ 是一个外测度（特别地，$\eta_\infty^*$ 即\eqref{align:ch01-外测度生成a} 中定义的外测度 $\mu^*$）。容易看出映射
\begin{align}
\eta^*:\mathcal{P}(X)\to[0,\infty],\qquad
\eta^*(E)\triangleq \lim_{\delta\to0^+}\eta_\delta^*(E)=\sup_{\delta>0}\eta_\delta^*(E)\quad(E\subseteq X) \label{align:ch01-eta星}
\end{align}
也是一个外测度。事实上，对任意 $E\subseteq X$，映射 $\delta\mapsto\eta_\delta^*(E)$ 非增；故定义是适定的，且
\begin{align}
\eta_\infty^*(E)\leqslant\eta_\delta^*(E)\leqslant\eta^*(E)\qquad\text{对任意 }\delta\in(0,\infty). \label{align:ch01-eta不等式}
\end{align}
于是对任意 $\{E_n\}\subseteq\mathcal{P}(X)$ 和 $\delta>0$，
\[
\eta_\delta^*\left(\bigcup_{n=1}^\infty E_n\right)\leqslant\sum_{n=1}^\infty\eta_\delta^*(E_n)\leqslant\sum_{n=1}^\infty\eta^*(E_n)
\quad\Rightarrow\quad
\eta^*\left(\bigcup_{n=1}^\infty E_n\right)\leqslant\sum_{n=1}^\infty\eta^*(E_n).
\]
这就证明了 $\eta^*$ 是 $\sigma$--次可加的。显然它是单调的，$\eta^*(\varnothing )=0$，故 $\eta^*$ 是一个外测度。以下引理完成该断言。

\begin{lemma}\label{lemma:ch01-eta星是度量外测度}
设 $(X,d)$ 为一度量空间。则由\eqref{align:ch01-eta星} 定义的外测度 $\eta^*$ 是一个度量外测度。
\end{lemma}

\begin{proof}
设 $E,F\subseteq X$ 非空，$d(E,F)>0$。若证得 $\eta^*(E\cup F)\geqslant\eta^*(E)+\eta^*(F)$，结论即成立。设 $\eta^*(E\cup F)<\infty$，否则结论显然。取 $\delta\in(0,d(E,F))$，设 $\{G_n\}\subseteq\mathcal{C}$ 且对所有 $n\in\mathbb{N}$ 有 $\operatorname{diam}(G_n)\leqslant\delta$，使得 $E\cup F\subseteq\bigcup_{n=1}^\infty G_n$。显然，由 $\delta$ 的选取，不存在 $\bar{n}\in\mathbb{N}$ 使得 $G_{\bar{n}}$ 同时与 $E$ 和 $F$ 相交。于是对所有 $n\in\mathbb{N}$，存在 $G_n',G_n''\subseteq X$ 使得 $G_n'\cup G_n''=G_n$；因而
\[
\max\{\operatorname{diam}(G_n'),\operatorname{diam}(G_n'')\}\leqslant\delta,
\]
且 $E\subseteq\bigcup_{n=1}^\infty G_n'$，$F\subseteq\bigcup_{n=1}^\infty G_n''$。显然这蕴含
\[
\eta_\delta^*(E)+\eta_\delta^*(F)\leqslant\eta_\delta^*(E\cup F),
\]
令 $\delta\to0^+$，结论成立。
\end{proof}

\begin{remark}\label{remark:ch01-eta是容量}
由\cref{remark:ch01-容量与外测度}，\eqref{align:ch01-eta-delta}--\eqref{align:ch01-eta星} 中定义的外测度 $\eta_\delta^*$（$\delta\in(0,\infty]$）和 $\eta^*$ 都是 $X$ 上的容量。
\end{remark}

度量外测度可刻画如下。

\begin{proposition}\label{proposition:ch01-度量外测度等价}
设 $(X,d)$ 为一度量空间，$\mu^*:\mathcal{P}(X)\to[0,\infty]$ 为一外测度。则以下陈述等价：
\begin{enumerate}[(i)]
\item 每个 Borel 集都是 $\mu^*$--可测的；
\item $\mu^*$ 是一个度量外测度。
\end{enumerate}
\end{proposition}

\begin{proof}
(i)$\Rightarrow$(ii)。由假设，对任意开集 $A\subseteq X$，有
\begin{align}
\mu^*(Z)=\mu^*(Z\cap A)+\mu^*(Z\cap A^c)\qquad\text{对任意 }Z\subseteq X \label{align:ch01-等价证明1}
\end{align}
（见\eqref{align:ch01-Caratheodory条件}）。设 $E,F\subseteq X$ 非空，$0<\delta<d(E,F)$。则 $G\triangleq \{x\in X\mid d(x,E)<\delta\}$ 是开集，且 $E\subseteq G$，$F\subseteq G^c$。于是在\eqref{align:ch01-等价证明1} 中取 $Z=E\cup F$、$A=G$，即化为\eqref{align:ch01-度量外测度条件}。结论成立。

(ii)$\Rightarrow$(i)。只需证每个闭集 $C\subseteq X$ 都是 $\mu^*$--可测的。对任意集合 $D\subseteq C^c$，令
\[
D_n\triangleq \{x\in D\mid d(x,C)\geqslant 1/n\}\subseteq D\qquad(n\in\mathbb{N}).
\]
现证
\begin{align}
\lim_{n\to\infty}\mu^*(D_n)=\mu^*(D). \label{align:ch01-等价证明2}
\end{align}
事实上，对所有 $n\in\mathbb{N}$，有 $\mu^*(D_n)\leqslant\mu^*(D)$，因为 $D_n\subseteq D_{n+1}\subseteq D$。还需证
\begin{align}
\lim_{n\to\infty}\mu^*(D_n)\geqslant\mu^*(D). \label{align:ch01-等价证明3}
\end{align}
设 $\lim_{n\to\infty}\mu^*(D_n)<\infty$（否则结论显然），令 $P_n\triangleq D_{n+1}\setminus D_n$（$n\in\mathbb{N}$）。若 $P_{2k}\neq\varnothing $（$k=1,\cdots ,n$），容易看出
\[
d\left(\bigcup_{k=1}^n P_{2k},P_{2n+2}\right)>0,
\]
故
\[
\mu^*\left(\bigcup_{k=1}^{n+1}P_{2k}\right)=\mu^*\left(\bigcup_{k=1}^n P_{2k}\right)+\mu^*(P_{2n+2}),
\]
因为 $\mu^*$ 是度量外测度。归纳论证得
\[
\mu^*\left(\bigcup_{k=1}^n P_{2k}\right)=\sum_{k=1}^n\mu^*(P_{2k})\qquad\text{对所有 }n\in\mathbb{N}
\]
（注意若某些 $P_{2k}=\varnothing $，上述等式也成立）。类似地可证
\[
\mu^*\left(\bigcup_{k=0}^n P_{2k+1}\right)=\sum_{k=0}^n\mu^*(P_{2k+1}).
\]
由于对所有 $n\in\mathbb{N}$ 有 $\bigcup_{k=1}^n P_k\subseteq D_n$，得
\begin{align}
\sum_{k=1}^\infty\mu^*(P_k)
&=\lim_{n\to\infty}\left(\sum_{k=0}^n\mu^*(P_{2k+1})+\sum_{k=1}^n\mu^*(P_{2k})\right) \label{align:ch01-等价证明4}\\
&=\lim_{n\to\infty}\left(\mu^*\left(\bigcup_{k=1}^n P_{2k}\right)+\mu^*\left(\bigcup_{k=0}^n P_{2k+1}\right)\right)
\leqslant 2\lim_{n\to\infty}\mu^*(D_n)<\infty. \nonumber
\end{align}
最后，由于 $C$ 是闭集，对所有 $n\in\mathbb{N}$ 有 $D=D_n\cup\left(\bigcup_{k=n}^\infty P_k\right)$，故
\begin{align}
\mu^*(D)\leqslant\mu^*(D_n)+\sum_{k=n}^\infty\mu^*(P_k)\qquad(n\in\mathbb{N}). \label{align:ch01-等价证明5}
\end{align}
由\eqref{align:ch01-等价证明4} 有 $\lim_{n\to\infty}\sum_{k=n}^\infty\mu^*(P_k)=0$，故在\eqref{align:ch01-等价证明5} 中令 $n\to\infty$，得\eqref{align:ch01-等价证明3}。于是\eqref{align:ch01-等价证明2} 成立。

现固定 $Z\subseteq X$，取 $D=Z\cap C^c$。对任意 $n\in\mathbb{N}$，有
\[
Z\supseteq(Z\cap C)\cup D_n,
\]
故由 $\mu^*$ 的单调性，
\[
\mu^*(Z)\geqslant\mu^*((Z\cap C)\cup D_n)=\mu^*(Z\cap C)+\mu^*(D_n)\qquad(n\in\mathbb{N}),
\]
因为 $d(Z\cap C,D_n)\geqslant 1/n>0$ 且 $\mu^*$ 是度量外测度。另一方面，由\eqref{align:ch01-等价证明2} 有
\[
\lim_{n\to\infty}\mu^*(D_n)=\mu^*(D)=\mu^*(Z\cap C^c),
\]
于是在上述不等式中令 $n\to\infty$，结论成立。
\end{proof}

最后我们补充两个关于度量空间上 Borel 测度的概念。

\begin{definition}\label{definition:ch01-加倍测度}
设 $(X,d)$ 为一度量空间。$X$ 上的正则 Borel 测度 $\mu$ 称为：
\begin{enumerate}[(i)]
\item 维数 $s>0$ 的，若存在 $c\geqslant 1$ 使得对所有 $x\in X$ 和 $r>0$ 有
\[
\frac{1}{c}r^s\leqslant\mu(B(x,r))\leqslant cr^s;
\]
\item 加倍的，若存在 $c_D\geqslant 1$ 使得对所有 $x\in X$ 和 $r>0$ 有
\[
\mu(B(x,2r))\leqslant c_D\mu(B(x,r)).
\]
\end{enumerate}
\end{definition}

\begin{remark}\label{remark:ch01-加倍性质}
(i) 若 $\mu$ 是维数 $s$ 的，则它是加倍的，常数为 $c_D=c^2 2^s$。

(ii) 若 $\mu$ 是加倍的，则容易看出对所有 $x\in X$ 和 $r>0$ 有 $\mu(B(x,r))<\infty$，且若 $\mu\not\equiv 0$，则 $\mu(B(x,r))>0$。

(iii) 若 $\mu$ 是加倍的且 $0<r<R$，则
\begin{align}
\mu(B(x,R))\leqslant c_D\left(\frac{R}{r}\right)^s\mu(B(x,r)) \label{align:ch01-加倍估计}
\end{align}
其中 $s=\log_2 c_D=\frac{\log c_D}{\log 2}$。为证\eqref{align:ch01-加倍估计}（若 $\mu\not\equiv 0$；否则\eqref{align:ch01-加倍估计} 显然），令 $k\triangleq \left[\log_2\left(\frac{R}{r}\right)\right]$ 为 $\log_2(R/r)$ 的整数部分。由于 $2^k\leqslant R/r\leqslant 2^{k+1}$，利用 $\mu$ 的单调性和迭代加倍条件，得
\[
\mu(B(x,R))\leqslant\mu(B(x,2\cdot 2^k r))\leqslant c_D\mu(B(x,2^k r))
\leqslant c_D^2\mu(B(x,2^{k-1}r))\leqslant\cdots\leqslant c_D^{k+1}\mu(B(x,r)).
\]
若 $c_D=1$，则上述不等式给出\eqref{align:ch01-加倍估计}，因为此时 $s=0$。若 $c_D>1$，则由
\[
\frac{\mu(B(x,R))}{\mu(B(x,r))}\leqslant c_D^{k+1},
\]
得
\[
\log_{c_D}\left(\frac{\mu(B(x,R))}{c_D\mu(B(x,r))}\right)\leqslant k\leqslant\log_2\left(\frac{R}{r}\right)
=\log_{c_D}\left(\frac{R}{r}\right)\log_2 c_D
=\log_{c_D}\left(\frac{R}{r}\right)^s,
\]
故在此情形下\eqref{align:ch01-加倍估计} 也成立。证毕。
\end{remark}




\end{document}